%==================================================================

\subsection{Множество ростков}
\leavevmode\par
\begin{defn}
Положим
\begin{equation}\label{eq:etale-set}
 \et(\Shv{F}):=\coprod_{x\in X}\Shv{F}_x
\end{equation}
--- разность множеств слоёв, и обозначим через
$\pi\colon\et(\Shv{F})\to X$ проекцию, отображающую слой $\Shv{F}_x$
в точку $x$. Тогда $\pi^{-1}(x)=\Shv{F}_x$.
\end{defn}

Пока это просто множество: точки пространства --- это пары
<<точка базы + росток в ней>>. Осталось дать топологию.

\subsection{Топология}
\leavevmode\par
\begin{defn}[топология на $\et(\Shv{F})$]
На $\et(\Shv{F})$ вводится наиболее тонкая топология, при которой все
отображения
\[
 s\colon U\longrightarrow\et(\Shv{F}),\qquad
 s(x)=s_x\ \ (s\in\Shv{F}(U)),
\]
непрерывны.
\end{defn}

Здесь стоит понять, что именно определено. У нас есть семейство
отображений $s\colon U\to\et(\Shv{F})$ (одно для каждого открытого $U$
и каждого $s\in\Shv{F}(U)$). Требуем, чтобы они все были непрерывными,
и берём наиболее тонкую топологию, которая это обеспечивает: множество
$G\subset\et(\Shv{F})$ объявляется открытым тогда и только тогда, когда
прообраз $s^{-1}(G)$ открыт для всякой такой секции $s$. Семейство таких
множеств образует топологию. Эквивалентная, <<вычислимая>> формулировка
--- критерий открытости.

\begin{thm}[критерий открытости]\label{thm:open}
Подмножество $G\subset\et(\Shv{F})$ открыто тогда и только тогда, когда
для всякого открытого $U\subset X$ и всякого $s\in\Shv{F}(U)$ множество
\begin{equation}\label{eq:crit}
 \{\,x\in U:\ s_x\in G\,\}=s^{-1}(G)\subset U
\end{equation}
открыто в $U$.
\end{thm}

\begin{proof}
($\Leftarrow$) Обозначим через $\tau$ семейство подмножеств $G$,
удовлетворяющих условию. Проверим, что $\tau$ --- топология.
$\varnothing$ и $\et(\Shv{F})$ проходят условие очевидно. Объединение:
если $G=\bigcup_\alpha G_\alpha$, то
$s^{-1}(G)=\bigcup_\alpha s^{-1}(G_\alpha)$ --- объединение открытых.
Пересечение: $s^{-1}(G_1\cap G_2)=s^{-1}(G_1)\cap s^{-1}(G_2)$ ---
пересечение открытых. Значит $\tau$ --- топология, и по построению
все $s$ в ней непрерывны. Любая другая топология, делающая все $s$
непрерывными, состоит только из множеств, удовлетворяющих условию
\eqref{eq:crit}, то есть содержится в $\tau$. Поэтому $\tau$ --- наиболее
тонкая такая топология.

($\Rightarrow$) Пусть топология на $\et(\Shv{F})$ выбрана так, как в
определении, и $G$ открыто. Тогда $s^{-1}(G)$ --- прообраз открытого
множества при непрерывном отображении, значит открыт.
\end{proof}

\sectionsrc{формулировка критерия --- photo\_2 (стр.~10) и
 photo\_3 (стр.~11): <<подмножество $G\subset\et(\Shv{F})$
 открыто $\Leftrightarrow$ множество всех ростков, в которые попадают
 локации сечения, открыто»; на стр.~11 отдельно проверяется, что
 прообраз $\pi^{-1}(U)$ открыт: <<$H=\pi^{-1}(U)$ состоит из
 всевозможных ростков $s_x$, $x\in U$, $s\in\Shv{F}(U)$, тогда оно
 открыто, т.к. $U$ открыто».}

\subsection{\texorpdfstring{$\pi$}{} --- локальный гомеоморфизм}
\leavevmode\par
\begin{thm}\label{thm:localhomeo}
Отображение $\pi\colon\et(\Shv{F})\to X$ является локальным
гомеоморфизмом: для каждой точки $\xi\in\et(\Shv{F})$ найдётся открытая
окрестность, на которой $\pi$ --- гомеоморфизм на открытое подмножество
$X$.
\end{thm}

\begin{proof}
Сначала проверим непрерывность $\pi$. Для любого открытого $O\subset X$
и всякой секции $t\in\Shv{F}(V)$ имеем
$t^{-1}(\pi^{-1}(O))=V\cap O$, открытое в $V$; по критерию
открытости $\pi^{-1}(O)$ открыто.
Пусть $\xi=s_x\in\Shv{F}_x$ для некоторого $s\in\Shv{F}(U)$.
Рассмотрим отображение $\widetilde s:U\to\et(\Shv{F})$,
$y\mapsto s_y$. Оно непрерывно по определению топологии и
$\pi\circ\widetilde s=\id_U$, поэтому инъективно. Покажем, что его
образ открыт. Для любой секции $t\in\Shv{F}(V)$ по критерию
\ref{thm:open} имеем
\[
 t^{-1}(\widetilde s(U))=\{y\in V\cap U:t_y=s_y\}.
\]
Если $t_y=s_y$, то по определению равенства ростков существует открытая
окрестность $W_y\subset V\cap U$ точки $y$, на которой
$t|_{W_y}=s|_{W_y}$. Значит, указанное множество открыто в $V$;
следовательно, $\widetilde s(U)$ открыто в $\et(\Shv{F})$.
Ограничение $\pi|_{\widetilde s(U)}$ непрерывно, а его обратное ---
$\widetilde s$, поэтому это гомеоморфизм на открытое $U$. Так как
$\xi\in\widetilde s(U)$, $\pi$ является локальным гомеоморфизмом.
\end{proof}

Следствие: $(\et(\Shv{F}),\pi)$ --- étale-пространство над $X$:
у каждой его точки есть окрестность, отображающаяся гомеоморфно на
открытое подмножество базы. В общем случае $\pi$ не обязано быть
накрытием и $\pi^{-1}(U)$ не обязано быть произведением
$U\times(\text{дискретное множество})$.

\begin{bookbox}
\booksrc{Годеман, гл.~II, \S\,1, п.~1.1, с.~129---130 --- язык
 расслоений. \emph{Расслоённым пространством} с базой $X$ называют
 пару $(E,p)$, где $E$ --- топологическое пространство, а
 $p\colon E\to X$ непрерывно; множество $E(x)=p^{-1}(x)$ обычно
 называют слоем над точкой $x$ (примечание ред., с.~129).
 \emph{Сечением} над подмножеством $M\subset X$ называется
 непрерывное $s\colon M\to E$ с $p\circ s=\id$; множество $E(U)$
 всех сечений над открытым $U$, с ограничениями, --- это пучок
 (пучок ростков сечений расслоённого пространства). Примечание
 переводчика там же: автор писал \emph{espace étalé}, русский
 перевод предпочёл <<накрывающее пространство>> --- в более
 широком, чем обычно, смысле; \emph{ét} в $\et(\Shv{F})$ ---
 от этого же термина.}
\end{bookbox}

\begin{center}
\normalsize
\setlength{\tabcolsep}{6pt}
\renewcommand{\arraystretch}{1.25}
\begin{tabular}{@{}>{\raggedright\arraybackslash}p{0.30\textwidth}
 >{\raggedright\arraybackslash}p{0.46\textwidth}
 >{\raggedright\arraybackslash}p{0.15\textwidth}@{}}
\toprule
\textbf{Из Годемана (расслоения)} &
\textbf{В этой лекции (пучки)} &
\textbf{Где} \\
\midrule
пара $(E,p)$, база $X$ & $\pi\colon\et(\Shv{F})\to X$ & \S\,4.1 \\
слой $E(x)=p^{-1}(x)$ & слой ростков $\Shv{F}_x$ & \S\,3.1 \\
сечение $s$, $p\circ s=\id$ & элемент $\Shv{F}^{\#}(U)$ &
 \S\,4.4--4.5 \\
$p$ --- локальный гомеоморфизм (накрывающее пространство) &
 $\pi$ --- локальный гомеоморфизм & \S\,4.3 \\
пучок ростков сечений $U\mapsto E(U)$ & $\Shv{F}^{\#}$ & \S\,5 \\
изоморфизм расслоений & изоморфизм пучков & \S\,8 \\
\bottomrule
\end{tabular}
\end{center}
\par\smallskip

\subsection{Сечение как непрерывный разрез}
Рисунок~\ref{fig:etale} показывает конструкцию целиком: над каждой
точкой $X$ висит свой слой, а сечение --- это <<разрез>>, который над
каждой точкой выбирает росток и делает это непрерывно.

\begin{risunok}{Проекция $\pi\colon\et(\Shv{F})\to X$: над каждой
 точкой --- слой ростков; сечение $s\in\Shv{F}^{\#}(U)$ --- непрерывный
 разрез, проецирующийся в себя ($\pi\circ s=\id$).\label{fig:etale}}
\begin{tikzpicture}[>=Stealth,font=\small]
  % база X с открытым U
  \draw[fig base] (0,0) ellipse (3.6 and 0.9);
  \node[below right] at (3.3,-0.55) {$X$};
  \begin{scope}
    \clip (0,0) ellipse (3.6 and 0.9);
    \fill[primary!10] (-1.8,0) ellipse (1.9 and 0.85);
  \end{scope}
  \draw[fig open] (-1.8,0) ellipse (1.9 and 0.85);
  \node[primary,below] at (-1.8,-1.25) {$U$};
  % этажёрка ét(F) — без рамки
  \node[anchor=north west] at (-4.15,5.15) {$\et(\Shv{F})$};
  % волокна и ростки
  \draw[fig fiber] (-2.6,0.63) -- (-2.6,4.9);
  \draw[fig fiber] (-0.9,0.88) -- (-0.9,4.9);
  \draw[fig fiber] ( 1.4,0.83) -- ( 1.4,4.9);
  \draw[fig fiber] ( 3.0,0.50) -- ( 3.0,4.9);
  \foreach \x/\y in {-2.6/1.7,-2.6/3.8,-0.9/2.2,-0.9/4.2,
                     1.4/1.9,1.4/3.9,3.0/2.3,3.0/4.3}
    \fill (\x,\y) circle (2.4pt);
  \foreach \x/\y in {-2.6/2.6,-0.9/3.2}
    \fill[danger] (\x,\y) circle (2.6pt);
  % сечение над U
  \draw[fig sect,smooth] plot coordinates
    {(-3.45,2.35) (-2.6,2.6) (-0.9,3.2) (0.15,3.05)};
  \node[danger] at (-0.3,3.62) {$s$};
  % слой над точкой x
  \draw[decorate,decoration={brace,mirror,raise=4pt}]
    (-2.6,1.7) -- (-2.6,3.8);
  \node[rotate=90,anchor=south,font=\footnotesize] at (-3.4,2.75)
    {$\pi^{-1}(x)=\Shv{F}_x$};
  % проекция
  \draw[->,thick] (4.75,5.0) -- (4.75,0.35);
  \node at (4.98,2.7) {$\pi$};
\end{tikzpicture}
\end{risunok}

Пояснение к рисунку~\ref{fig:etale}. Прямоугольник сверху ---
пространство ростков $\et(\Shv{F})$, овал снизу --- база $X$.
Вертикальные пунктирные линии --- волокна $\pi^{-1}(x)$: над
каждой точкой $x$ висит свой слой ростков (столбик чёрных точек),
скобкой слева показано волокно над одной точкой,
$\pi^{-1}(x)=\Shv{F}_x$. Синим выделено открытое $U\subset X$;
красная кривая $s$ --- сечение над $U$: для каждой точки $x\in U$
она выбирает росток (красные точки на волокнах), причём выбор
меняется непрерывно, поэтому кривая не рвётся. Стрелка $\pi$
--- проекция вниз, и $\pi\circ s=\id$ означает: каждый выбранный
росток лежит ровно над своей точкой.

\subsection{Сечение и мульти-секция}
Непрерывность $s$ означает, что выбор ростков по точкам $x\in U$
происходит <<плавно>>: множество $\{s(x):x\in U\}$ открыто в
$\et(\Shv{F})$. Полезно противопоставить этому произвольный выбор
ростков (рисунок~\ref{fig:multisection}): можно для каждой точки $U$
произвольно вытащить из слоя какой угодно росток --- получится
<<мульти-секция>>, но, вообще говоря, не сечение.

\begin{risunok}{Сечение --- непрерывный выбор ростков (красное); синий
 произвольный выбор ростков по точкам $U$ сечением не является.
 Формально $\Shv{F}^{\#}(U)=\{s\mid(\pi\circ s)(U)=\id_U\}$.
  \label{fig:multisection}}
\begin{tikzpicture}[>=Stealth,font=\small]
  \draw[fig base] (0,0) ellipse (3.4 and 0.75);
  \node[below] at (0,-0.8) {$U$};
  \foreach \x/\top in {-2.6/0.483,-1.3/0.693,0/0.75,1.3/0.693,2.6/0.483}
    \draw[fig fiber] (\x,\top) -- (\x,3.5);
  \foreach \x in {-2.6,-1.3,0,1.3,2.6}
    \foreach \y in {1.5,2.4,3.2} \fill (\x,\y) circle (2.4pt);
  % непрерывное сечение
  \draw[fig sect,smooth] plot coordinates
    {(-3.0,1.2) (-2.6,1.5) (-1.3,2.4) (0,3.2) (1.3,2.4) (2.6,1.5) (3.0,1.2)};
  \node[danger,above] at (0,3.55) {$s\in\Shv{F}^{\#}(U)$};
  % произвольная мульти-секция
  \fill[warning] (-2.6,2.4) circle (3.2pt);
  \fill[warning] (0,1.5) circle (3.2pt);
  \fill[warning] (2.6,3.2) circle (3.2pt);
  \draw[warning,dashed] (-2.6,2.4) -- (0,1.5) -- (2.6,3.2);
  \node[warning,font=\footnotesize] at (0,-1.95)
    {произвольный выбор ростков: мульти-секция, не сечение};
\end{tikzpicture}
\end{risunok}

Пояснение к рисунку~\ref{fig:multisection}. Внизу --- открытое
$U\subset X$, вверх от него идут волокна: столбики ростков над
точками $U$. Красная кривая --- сечение: она соединяет ростки в
непрерывный <<разрез>> и проецируется в $U$ «на себя». Синие
точки --- произвольный выбор ростка над каждой точкой $U$:
соединённые пунктиром, они скачут между уровнями, выбор
не непрерывен, поэтому сечением не является.

\begin{bookbox}
\booksrc{Годеман, гл.~II, \S\,1, п.~1.3, с.~133 --- сечения
 определяются над произвольным подмножеством $M\subset X$, не
 только над открытым. При этом (F1) остаётся осмысленной и верной
 без открытости, а (F2) --- нет: иначе любое отображение
 $s\colon X\to E$ с $p\circ s=\id$ было бы непрерывным, ведь
 $X$ распадается на одноточечные множества. Теорема~1.3.1 там же:
 склейка сечений работает на локально конечном \emph{замкнутом}
 покрытии.\par\smallskip
 Манин, \S\,1.10, с.~63 --- расслоенное произведение;
 \S\,1.12, с.~77---79 (определение~1.12.9, с.~79): векторное
 расслоение --- это семейство векторных пространств, локально
 тривиальное в каждой точке. Это прямой прототип будущего пучка
 модулей: <<расслоение>> в алгебраизированном виде.}
\end{bookbox}

\subsection{Нехаусдорфовость}
\leavevmode\par
\begin{rem}[нехаусдорфовость]
Даже если $X$ хаусдорфово, $\et(\Shv{F})$, вообще говоря, не
хаусдорфово. По Годеману (Замечание 1.2.4, с.~133): если
$\et(\Shv{F})$ отделимо, то два сечения $s,t$ над открытым $U$,
совпадающие в некоторой точке, равны во всей связной компоненте,
т.\,е. в пучке действует <<принцип аналитического продолжения>>.
Например, пучок ростков голоморфных функций на комплексном
многообразии отделим, а пучок ростков непрерывных функций --- нет.
\end{rem}

\begin{risunok}{Пример нехаусдорфова $\et(\Shv{F})$: $f_1$ и $f_2$
 совпадают на $[0,\tfrac12]$ и расходятся дальше. Ростки над точкой
 $\tfrac12$ различны, но любые их окрестности пересекаются, так как
 листы соединяются слева (иллюстрация к Манину, рис.~2.3, и к заметке
 Годемана 1.2.4).\label{fig:nonhausdorff}}
\begin{tikzpicture}[>=Stealth,font=\small]
  % левый график: f_1 и f_2
  \draw[->] (-0.3,0) -- (4.45,0);
  \draw[->] (0,-0.3) -- (0,3.15);
  \draw[fig sect,dashed] (0.3,0.15) -- (4.15,0.15);
  \draw[fig leafA,smooth] plot coordinates
    {(0.3,0.15) (1.7,0.15) (2.5,0.5) (3.3,1.4) (4.1,2.45)};
  \draw (1.7,0) -- (1.7,-0.09);
  \node[below] at (1.7,-0.36) {$\tfrac12$};
  \node[below left] at (0.05,-0.1) {$0$};
  \node[primary] at (3.95,2.62) {$f_1$};
  \node[danger] at (3.5,0.45) {$f_2$};
  % правая панель: ét(F)
  \begin{scope}[xshift=6.4cm]
    \node[anchor=south west] at (-0.3,3.05) {$\et(\Shv{F})$};
    \draw[line width=1.4pt] (-0.3,2.1) -- (1.7,2.1);
    \draw[fig leafA,smooth] plot coordinates
      {(1.7,2.1) (2.7,2.45) (4.3,2.85)};
    \draw[fig sect,smooth] plot coordinates
      {(1.7,2.1) (2.7,1.65) (4.3,1.25)};
    \fill (2.7,2.45) circle (2.6pt);
    \fill (2.7,1.65) circle (2.6pt);
    \draw[<->,gray!70!black,line width=0.6pt] (2.7,1.82) -- (2.7,2.28);
    \node[fig note,align=center] at (3.35,0.55)
      {окрестности\\пересекаются};
    \node[primary] at (4.0,2.45) {$f_1$};
    \node[danger] at (3.7,1.25) {$f_2$};
    \draw[fig fiber] (1.7,2.1) -- (1.7,0.15);
    \node at (1.7,-0.36) {$\tfrac12$};
  \end{scope}
\end{tikzpicture}
\end{risunok}

Пояснение к рисунку~\ref{fig:nonhausdorff}. Слева --- графики
двух функций: синий $f_1$ и красный пунктирный $f_2$ совпадают
на отрезке $[0,\tfrac12]$ (участок одного цвета) и дальше
расходятся. Справа --- как это выглядит в $\et(\Shv{F})$: слева
от отметки $\tfrac12$ обе кривые лежат в одном и том же листе
(общий чёрный отрезок), а справа расходятся на два разных листа.
Над точкой $\tfrac12$ два ростка --- две разные точки (чёрные
точки, соединённые стрелкой), но любые их окрестности
обязательно пересекаются, раз листы слева соединены. Отсюда:
$\et(\Shv{F})$ не хаусдорфово.

\begin{bookbox}
\booksrc{Годеман, гл.~II, \S\,1, п.~1.2, с.~130---133: топология
  вводится как <<наиболее тонкая, при которой отображения $s$
 непрерывны>>, и доказывается критерий открытости.
 \textbf{Теорема 1.2.1 (Годеман, с.~132):} всякий пучок множеств с
 базой $X$ изоморфен пучку ростков некоторого накрывающего пространства
 над $X$, определённого однозначно с точностью до изоморфизма.
\par\smallskip
 Манин, \S\,2.1.8, с.~113---114 --- эквивалентное определение пучка:
 пара $(Y_{\Shv{F}},r\colon Y_{\Shv{F}}\to X)$, где $r$ локально ---
 гомеоморфизм, а сечения $\Shv{F}(U)$ --- отображения $s\colon U\to
 Y_{\Shv{F}}$ с $r\circ s=\id$. Обратно, из пучка получается
 $Y_{\Shv{F}}=\coprod_x\Shv{F}_x$ с топологией, порождённой сечениями.
}
\end{bookbox}

%==================================================================
